\subsection{Background and related work}

Every finite pattern is measure universal. Indeed, if $F$ is finite
and $A$ is measurable with finite positive measure, continuity of
translations in $L^1$ gives
\[
 m\Bigl(\bigcap_{a\in F}(A-sa)\Bigr)\longrightarrow m(A)>0
 \qquad(s\to0).
\]
For an arbitrary positive-measure set, apply this observation to a
bounded subset of positive measure. The conjecture asks whether this
property always fails for infinite patterns.

Falconer~\cite{Falconer1984} and Eigen~\cite{Eigen1985} proved
nonuniversality for positive decreasing sequences $a_n\to0$ with
$a_{n+1}/a_n\to1$. Humke and Laczkovich
\cite[Theorems~1--2 and Corollary~1]{HumkeLaczkovich1998} developed a
finite covering characterization and further criteria in terms of
relative gaps. Kolountzakis~\cite[Theorem~3]{Kolountzakis1997} gave
a probabilistic criterion using large finite subsets with sufficiently
large normalized minimum gap. In the translation-invariant form of
Chleb\'ik~\cite[Theorem~15]{Chlebik2015}, a bounded infinite set is
nonuniversal if it contains finite subsets $a_1>\cdots>a_m$, of
arbitrarily large cardinality, for which
\[
 -\log\left(\frac{\min_{i<m}(a_i-a_{i+1})}{a_1-a_m}\right)=o(m).
\]
Geometric progressions do not satisfy this criterion. For distinct
powers of $q$, we have $a_{m-1}\le q^{m-2}a_1$ and
$a_1-a_m\ge(1-q)a_1$. Their normalized minimum gap is therefore at
most $q^{m-2}/(1-q)$, whose negative logarithm cannot be sublinear
in $m$.

Additive structure supplies another family of results.
Bourgain~\cite{Bourgain1987} proved that the sum of three infinite
sets is nonuniversal. Kolountzakis~\cite{Kolountzakis1997} treated
certain double sums, including $G_{1/2}+G_{1/2}$.
Mora Cuellar, Iosevich, Kulkarni, Rojas Aravena and
Yavicoli~\cite[Theorem~2]{MoraCuellarEtAl2026Twofold} proved that the
sum or difference of a geometric null sequence and any infinite set
is nonuniversal. Avoiding such a larger pattern does not imply
avoidance of the single geometric progression in Theorem~\ref{thm:main}.

Other results concern different classes of patterns or notions of
largeness. Iosevich, Kulkarni, Mora Cu\'ellar, Rojas Aravena and
Yavicoli~\cite[Theorem~1.5]{IKMRY2026Rajchman} prove nonuniversality
for sets supporting a probability measure whose Fourier transform
tends to zero at infinity. Such a measure is atomless, so the
hypothesis excludes countable sets
\cite[Proposition~4.1]{IKMRY2026Rajchman}.
Cruz, Lai and Pramanik~\cite[preprint, Corollary~1.2 and
Proposition~4.4]{CruzLaiPramanik2023} construct sets of Hausdorff
dimension one avoiding certain null sequences together with their
limit point; the subsets establishing that dimension have Lebesgue
measure zero. In a different direction, the theorem of Feng, Lai and
Xiong~\cite[preprint, Theorem~1.1]{FengLaiXiong2024} implies that
every $G_q$ embeds into every measurable set of positive measure by
a bi-Lipschitz map. The restriction to affine maps is therefore
substantive.

The proof below belongs to the probabilistic approach of
Kolountzakis~\cite{Kolountzakis1997}. Random cells, fixed-center
scale discretization, and integration of exceptional-center
probabilities also appear in Chleb\'ik~\cite[Section~5]{Chlebik2015}
and Kolountzakis--Papageorgiou~\cite[Section~3.1]{KolountzakisPapageorgiou2025};
the latter paper treats unbounded sequences. Solymosi~\cite[Section~3]{Solymosi2011Sparse}
and Tom~\cite{Tom2015} report incomplete dyadic constructions based on
random cells; Tom also considers nested dyadic intervals and tests
independent of the dilation. Our contribution is the
finite routing construction and its local control of the scale
count. We explain that mechanism next and give a self-contained proof.
